
\textbf{迭代器}，是一种检查容器内元素并遍历元素的数据类型，通常用于对C++中各种容器内元素的访问，
但不同的容器有不同的迭代器，初学者可以将迭代器理解为一个可以\textbf{自动向后移动的指针}。

常用的可迭代的容器有：

\begin{enumerate}
    \item 序列容器：$array$、$vector$、$deque$、$list$、$forward\_list$等
    \item 关联容器：$set$、$map$、$multimap$等
\end{enumerate}

$queue$、$stack$、$priority\_queue$是容器适配器，必须满足一定规则，没有迭代器。

对于可迭代容器，一般可以使用.begin()和.end()方法获取迭代器，.begin()方法返回指向第一个元素的迭代器，
.end()方法返回指向最\textbf{后一个元素的迭代器的下一位置}。
使用迭代器，可以方便的对容器内元素进行遍历，例如：

\begin{tcode}
std::vector<int> vec{1,2,3,4,5,6,7,8,9,0};
for(auto it = vec.begin();it != vec.end();++it){
    std::cout << *it << std::endl;
}
\end{tcode}

也可以简写为：

\begin{tcode}
for (auto it: vec) {
    std::cout << it << std::endl;
}
\end{tcode}

还可以使用迭代器对容器内元素进行修改，例如：

\begin{tcode}
std::vector<int> vec{1, 2, 3, 4, 5, 6, 7, 8, 9, 0};
vec.erase(vec.begin()+5);//删除第6个元素
vec.erase(vec.begin()+5,vec.end());//删除第6个元素之后的所有

std::vector<int> vec1{1, 2, 3, 4, 5, 6, 7, 8, 9, 0};
std::iter_swap(vec1.begin(),vec1.begin()+1);//交换第一个和第二个元素

std::vector<int> vec2{1, 2, 3, 4, 5, 6, 7, 8, 9, 0};
//此函数需要头文件<algorithm>
std::sort(vec2.begin(),vec2.end(),[](int& a,int& b){
    return a>=b;
});

//此函数需要头文件<algorithm>
//删除大于5的元素
auto newVec = std::remove_if(vec2.begin(),vec2.end(),[](int& a){
    return a>5;
});
\end{tcode}

我们不会使用十分复杂的迭代器，感兴趣的同学可以看一下cppreference的相关文档：

\url{https://zh.cppreference.com/w/cpp/iterator}

\textbf{迭代器失效}：迭代器失效指的是迭代器因为某些操作而不再指向原来的元素，
或者不再保持有效的状态，这样的迭代器如果继续使用，可能会导致未定义行为，包括程序崩溃。
通俗的理解就是，当你对容器进行了操作时，更改了容器的大小，但是迭代器仍然指向原来的位置，
这样就会导致读取错误的内存，这就是迭代器失效。

\begin{tcode}
std::vector<int> vec={1,2,3,4,5,6,7,8,9,10};
for(auto it = vec.begin(); it!= vec.end(); ++it){
    if(*it%2==0){
        vec.erase(it);
    }
}
\end{tcode}

看上面的代码，当执行到vec.erase(it)时，it已经失效，因为执行了erase之后，
vec的大小发生了变化，这导致了迭代器失效。

在实践过程中， 在下面的几种情况下，迭代器会失效，需要\textbf{避免}出现：
\begin{enumerate}
    \item 当执行erase方法时，指向删除节点的迭代器全部失效，指向删除节点之后的全部迭代器也失效
    \item 当进行push\_back\(\)方法时，end操作返回的迭代器肯定失效。
    \item 当插入\(push_back\)一个元素后，capacity返回值与没有插入元素之前相比有改变，
    则需要重新加载整个容器，此时first和end操作返回的迭代器都会失效。
    \item 当插入\(push_back\)一个元素后，如果空间未重新分配，指向插入位置之前的元素的迭代器仍然有效，但指向插入位置之后元素的迭代器全部失效。
\end{enumerate}

对于关联容器，使用迭代器时，删除一个键值对知识当前这个迭代器失效，别的迭代器不会受到影响，
这可能是由于底层使用红黑树实现的原因，但是实际应用中，仍然不应该冒这个风险，因此，
必须\textbf{坚决避免迭代器失效问题的发生}。

在具体的应用中，我们在使用迭代器遍历时，应该注意不要对容器进行增删操作，更不应该改变容器的大小。
可以将不需要的元素放在\textbf{容器的末尾}或者\textbf{记录下需要删除的迭代器}，然后适时删除。例如：

\begin{tcode}
for(auto it = vec.begin(); it!= vec.end(); ++it){
    if((*it)%2==0){
        its.push_back(it-vec.begin());
    }
}

int i=0;
for(auto it : its){
    vec.erase(vec.begin()+it-i);
    i++;
}
}
\end{tcode}

这种方法，非常\textbf{不优雅}，因此，\textbf{不要使用}。

\textbf{更多的用法等待着同学们自己去探索！}

